\chapter{Introduction}
\label{ch:introduction}
\section{Background and Problem}
Curved channels couple fluid turning with particle inertia. A fluid marker
follows the local velocity, whereas a finite-inertia particle can lag behind
a change in flow direction and approach a wall. Predicting the resulting
trajectory is relevant to transport, separation, and deposition studies.
The mathematical difficulty is broader than approximating a smooth velocity
field: small, spatially localized errors can change a discrete outcome such
as whether a particle exits or sticks.

Physics-informed neural networks incorporate governing-equation and boundary
residuals into a learning objective \cite{raissi2019,karniadakis2021}. They
offer a differentiable approximation that can be evaluated during trajectory
integration. Their optimization, scaling, and accuracy nevertheless require
verification; a small loss or an attractive plot is insufficient evidence
of a reliable solution \cite{wang2021,wang2022}. Classical CFD remains a
necessary independent comparator in the present study.

The current project already includes a two-network implementation, analytical
tests, an independent FEM reference, and particle tracking. Its preliminary
results reveal a useful research problem: the principal flow approximation
has a relatively small mean error, yet some particle paths and endpoint
classes differ from the reference. This proposal therefore studies the
connection between field accuracy and transport accuracy, rather than
assuming one implies the other.

\section{Aim and Objectives}
The aim is to develop and assess a reproducible physics-informed workflow
for predicting dilute-particle transport in smooth curved channels within
an explicitly defined physical and geometric scope.
\begin{enumerate}[label=O\arabic*.]
\item Quantify flow and pressure accuracy against analytical and independently
assembled numerical references, including mesh and interpolation sensitivity.
\item Determine how local field errors, inertia, particle size, and wall
clearance affect trajectory errors and endpoint classification.
\item Evaluate targeted changes to training and sampling using separate
development and final evaluation cases.
\item Establish the supported scope of geometry variation and a distinct
three-dimensional verification pathway, without inferring universal
generalization from the network input dimensions.
\end{enumerate}

\section{Research Questions}
\begin{enumerate}[label=RQ\arabic*.]
\item Under which training and reference-resolution conditions does the
flow PINN satisfy both field-error and mass-conservation criteria?
\item How does flow-approximation error propagate into particle position,
residence time, and sticking or exit decisions?
\item Does a learned force surrogate materially change trajectories compared
with exact Stokes drag evaluated on the same learned flow?
\item Which conclusions persist across separately trained channel shapes and
Reynolds numbers, and what further evidence is needed for a 3D channel model?
\end{enumerate}
The working expectation is that near-wall and turning-region accuracy will
be more informative for deposition than a global sampled velocity norm.
This is a testable hypothesis for the proposed study, not a mechanism
already established by the preliminary figures.

\section{Scope and Intended Contribution}
The initial physical model is steady, incompressible, Newtonian flow with
one-way coupling to dilute, heavy, spherical particles. Stokes drag and
irreversible wall sticking are retained. The supported channel family has
a smooth, nonintersecting planar centerline and constant width. Sharp
corners, junctions, variable-width networks, particle--particle forces,
and feedback on the fluid are outside the present implementation.

The intended contribution is a reproducible evaluation linking field errors
to transport errors. Its scope excludes a universal drag law and arbitrary
3D geometry. The remaining chapters develop the foundations, methods,
evaluation, preliminary evidence and proposed work.
